\section{Seasonal-window and screening-rule diagnostics}
The previously selected source sample contains 96 eligible long-record stations; this calculation uses the 95-station project-processed subset with a documented candidate redistribution path. None of the 35 recorded witnesses is at the excluded LUBW station, although it remains part of the original eligible-station denominator. The subset comprises \SWComplete{} complete station-years at \SWStations{} stations and \SWGroups{} connected groups. For each complete year, we compare the full-year reported H30 with the largest 30-day mean whose every day lies from 1 July through 31 August. Of those years, \SWAbove{} exceed the summer-window value by more than the numerical tolerance $5\times10^{-6}$ $^{\circ}$C, and \SWAboveHalf{} exceed it by more than 0.5 $^{\circ}$C. The mean difference is \SWMean{} $^{\circ}$C. These are complete recorded years, not simulated reconstruction errors; neither the tolerance nor the 0.5 $^{\circ}$C threshold is an estimate of sensor error.

The comparison was specified after the original witnesses and literature had been examined. It is not a reproduction of Beaufort et al.'s multiyear thermal peak or their summer air-temperature climate correction; that study also excluded sites with strong human influence. The present sample contains regulated and incompletely characterized sites, covers different locations and years, and the fully-contained July--August window differs from a peak-center criterion. Consequently our fraction must not be compared directly with their reported fraction. The incomplete-year comparison is also conditional on the 35 already selected record witnesses: all have a fully observed July--August 30-day window, and \SWWitness{} have a larger fully observed 30-day window elsewhere in the year. This does not establish the full H30 of an incomplete year or a population prevalence.

Table~\ref{tab:summer-rules} counts previously eligible station-years rejected by four explicit study rules, including how many still supply at least one fully observed 30-day window. A lower bound is information about the reported record, not a complete annual estimate or a guarantee of better trends. The 90\% station-specific warm-season rule uses a leave-target-year-out climatology; the complete July--August row is a contextual comparator rather than the full published climate-correction procedure. The records are not an independent validation sample and no sensor error has been estimated.

\begin{table}[ht]\caption{Study screening rules on \SWRuleYears{} eligible station-years in the 95-station candidate subset. Rejected-with-bound means the rejected incomplete year contains at least one complete observed 30-day window.}\label{tab:summer-rules}\centering\small
\resizebox{\linewidth}{!}{\begin{tabular}{lrrrrr}\toprule
Rule & Accepted & Rejected & Unknown & Rejected with bound & Stations \\
\midrule
Complete year & 2,122 & 534 & 0 & 518 & 70 \\
Annual 95\% & 2,550 & 106 & 0 & 90 & 48 \\
Station warm season 90\% & 2,590 & 66 & 0 & 50 & 30 \\
Complete July--August & 2,472 & 184 & 0 & 168 & 58 \\
\bottomrule\end{tabular}}\end{table}

After inspecting the largest annual gaps, we separately summarized existing project regulation labels and removed the two sites with the largest individual gaps (Table~\ref{tab:summer-influence}). The 2010 USGS analysis identified Blue River and Fall Creek as dam-affected monitoring sites in 2001--2002; it does not establish the cause or measurement quality of every year examined here. The existing \emph{unmatched gages} label does not mean unregulated, and labels and sites were not randomized. Removing these two post-hoc sites leaves \SWExclAbove{} of \SWExclYears{} complete station-years above the numerical tolerance, with mean difference \SWExclMean{} $^{\circ}$C. These summaries do not identify causal dam effects or transfer to streams absent from the sample.

\begin{table}[ht]\caption{Post-hoc seasonal-window influence by existing project label and after removing two selected sites. ``Above'' is greater than $5\times10^{-6}$ $^{\circ}$C; the second count exceeds 0.5 $^{\circ}$C. Labels are observational metadata, not independently verified treatments.}\label{tab:summer-influence}\centering\small
\resizebox{\linewidth}{!}{\begin{tabular}{lrrrrr}\toprule
Cohort & Stations & Years & Above & $>0.5$ & Mean ($^{\circ}$C) \\
\midrule
All stations & 95 & 2,122 & 779 & 304 & 0.255 \\
Exclude two selected sites & 93 & 2,086 & 749 & 279 & 0.197 \\
Label: regulated & 23 & 397 & 173 & 103 & 0.725 \\
Label: unmatched gages & 61 & 1,542 & 569 & 196 & 0.158 \\
Label: unknown & 11 & 183 & 37 & 5 & 0.049 \\
\bottomrule\end{tabular}}\end{table}

The comparison and screening totals were independently reconstructed using pandas rolling windows on the same snapshot. The post-hoc influence totals and 36 selected-site peak-window dates were independently checked using NumPy convolution; the stored regulation labels were checked against the project's source metadata. Neither audit verifies the original provider's measurements or whether each project regulation label is accurate. The supplied 95-station snapshot permits the window and screening arithmetic to be rerun, subject to its attribution and reuse conditions.
